% Chapter 2: Background Study and Literature Review
\chapter{Background Study and Literature Review}

\section{Introduction}

Suicide is shaped by psychological, social, economic, cultural, and environmental factors. In low- and middle-income countries (LMICs) such as Bangladesh, the challenge is compounded by the absence of a national suicide registry and pervasive under-reporting, which makes reliable, structured surveillance data extremely difficult to obtain. The most comprehensive case-level source available to researchers in this context remains newspaper and media reports --- a source rich in detail yet severely constrained by incompleteness and reporting bias \cite{ref1,ref50}. The chapter that follows reviews the current empirical picture of suicide in Bangladesh, the documented limitations of newspaper-derived datasets, established approaches to missing-data imputation in public health, and prior applications of clustering to suicide and mental-health research, before consolidating the gap that motivates this study.

\section{Suicide in Bangladesh: Current Knowledge}

National health surveys and WHO reports place the annual figure at approximately 10,000 or more deaths in Bangladesh, while more recent analyses of Bangladesh police records report a substantially higher average of roughly 14,719 suicides per year over 2020--2024 \cite{ref1,ref50,ref3}; both estimates are widely understood to be lower bounds, given pervasive under-reporting. A comparative analysis of the Bangladesh Health and Injury Surveys from 2002 and 2015 found suicide rates rising from 6.2 to 7.7 per 100,000 over that period, with the highest-risk group shifting from the elderly in 2002 to adolescents by 2015 \cite{ref26}.

The epidemiology shows several consistent patterns. Hanging accounts for approximately 65--82\% of completed suicides across studies \cite{ref50,ref51}, followed by poisoning, which is especially prevalent in rural areas where pesticides are readily available \cite{ref50}. Females --- particularly young married women and female students --- are disproportionately represented, reversing the global trend of male predominance \cite{ref50}. Meta-analytic and review evidence reports female predominance among reported cases, a substantial share of student and young-adult victims, and marital or familial discord as the most frequently cited risk factor \cite{ref50,ref51}. Rural areas show suicide rates an order of magnitude greater than urban centres \cite{ref50}, with specific districts --- Jhenaidah, Chuadanga, Meherpur, Narshingdi --- identified as high-burden. A dedicated study of Jhenaidah recorded a suicide rate of 20.6 per 100,000, well above national and global averages, with females showing higher rates than males across 2010--2018 \cite{ref29}.

Among vulnerable sub-populations, students have received growing attention. A retrospective analysis of online news portals identified 984 student suicide cases between 2022 and 2023, with females accounting for 61\% \cite{ref30}; contributing factors included academic pressure, board examination timing, family conflict, and failed relationships. COVID-19 further altered suicide patterns, with one study documenting a shift towards higher male incidence during the pandemic, attributed to economic disruption and job loss among primary breadwinners \cite{ref31}.

Despite this growing body of work, the majority of existing Bangladeshi studies share a critical limitation: they rely on manual or semi-automated extraction from newspaper sources and explicitly acknowledge the problem of incomplete records. Variables such as age, exact location, reason for suicide, socioeconomic status, and prior mental health history are frequently absent or inconsistently recorded, severely restricting the depth of analysis \cite{ref51,ref1}.

\section{Challenges with Newspaper-Based Suicide Datasets}

Newspaper and media-based datasets are the primary --- and often only --- source of case-level suicide data in resource-constrained settings such as Bangladesh, and a systematic review of media-based suicide research in LMICs concluded that content analyses of newspaper reports are the dominant research methodology in these settings \cite{ref23}.

This reliance carries well-documented structural limitations, the most critical being missing data. Reporters covering suicide incidents rarely record the full demographic, socioeconomic, and contextual detail required for meaningful quantitative analysis, so critical variables --- age, exact location, reason for suicide, socioeconomic indicators, environmental conditions, and prior mental health history --- are systematically absent across a large proportion of cases, with the prior-mental-health and family-structure fields effectively unrecoverable in practice \cite{ref51,ref1,ref50,ref29}. Per-feature missingness measured directly in the present study --- low-single-digit rates for fields such as gender and method through to over 90\% for previous mental health history and family size --- is reported quantitatively in Chapter~3 (Table~\ref{tbl:missingness_groups}) and is broadly consistent with the qualitative patterns described in the cited Bangladeshi literature.

Beyond missingness, newspaper-based datasets are inherently biased towards newsworthy cases --- unusual methods, public locations, prominent individuals --- so common, unremarkable suicides in remote rural areas are systematically under-reported \cite{ref23}. Reporting inconsistencies in spelling, categorical labels, and temporal formats further complicate preparation, and these challenges typically force researchers either into complete-case analysis (which sharply reduces sample size and introduces selection bias) or into naive imputation strategies such as mean substitution (which introduces distributional bias) \cite{ref6,ref55}.

\section{Related Work on Missing Data Imputation in Public Health}

Three families of imputation approaches dominate the recent public-health and clinical literature: Multiple Imputation by Chained Equations (MICE), K-Nearest Neighbours (KNN) imputation, and ensemble tree-based methods, most notably MissForest. MICE, formalised by Van Buuren and Groothuis-Oudshoorn (2011) \cite{ref10}, fits a series of regression models --- one per variable --- iteratively conditioned on all others; its flexibility for mixed-type data and principled treatment of uncertainty have made it widely recommended in clinical and epidemiological work. KNN imputation \cite{ref11} fills missing values from the weighted average of the $k$ most similar complete observations, is computationally efficient, and imposes no distributional assumptions, though its performance degrades in high dimensions or when local neighbourhood structure is disrupted by heavy missingness. MissForest \cite{ref12} trains a separate random forest per variable using all others as predictors and iterates until convergence; its non-parametric, ensemble-based design is particularly robust to non-linear relationships, interactions, and mixed types, and multiple comparative studies confirm its superiority over MICE and KNN on medical and biological datasets with moderate-to-high missingness and complex variable interactions \cite{ref12,ref32,ref33}.

Table~\ref{tbl:summary_of_related_work_on_mis} summarises the most relevant comparative imputation studies.

\begin{table}[H]
\centering
\caption[Related work on missing data imputation in public health]{Summary of related work on missing data imputation in public health and related domains}
\label{tbl:summary_of_related_work_on_mis}
\renewcommand{\arraystretch}{1.0}
\small
\begin{tabularx}{\textwidth}{|>{\hsize=0.790\hsize\linewidth=\hsize\raggedright\arraybackslash}X|>{\hsize=0.795\hsize\linewidth=\hsize\raggedright\arraybackslash}X|>{\hsize=0.824\hsize\linewidth=\hsize\raggedright\arraybackslash}X|>{\hsize=1.592\hsize\linewidth=\hsize\raggedright\arraybackslash}X|}
\hline
\textbf{Study} & \textbf{Domain} & \textbf{Methods Compared} & \textbf{Key Finding} \\ \hline
Stekhoven \& Bühlmann (2012) \cite{ref12} & Bioinformatics / mixed-type data & MissForest vs. MICE vs. KNN & MissForest outperforms MICE and KNN on mixed-type data, reducing imputation error by up to 60\% in some datasets. \\ \hline
Prakash et al. (2025) \cite{ref55} & Large-scale mental health surveys ({SPARK} autism cohort, 117k participants) & MICE, KNN, MissForest, MIDAS & MIDAS and KNN performed best on random and blockwise missingness; MICE was the most computationally costly but remained competitive at moderate missing rates. \\ \hline
Alkhateeb et al. (2024) \cite{ref32} & Medical datasets (10 datasets) & MICE, KNN, MissForest, RFE-MF & MissForest consistently outperforms KNN and MICE across diverse medical datasets at missingness rates of 10--50\%. \\ \hline
Hamdan et al. (2025) \cite{ref33} & Healthcare diagnostics & MICE, KNN, MissForest, mean/median & MissForest achieves lowest RMSE and MAE; MICE ranks second; KNN performance is inconsistent across datasets. \\ \hline
Jarrett et al. (2021) \cite{ref8} & 69 real-world tabular datasets & Multiple imputation methods & No single imputation method dominates across all data types and missingness patterns; context-specific evaluation is essential. \\ \hline
\end{tabularx}
\end{table}

Despite this growing body of comparative work, several important gaps remain. First, the vast majority of imputation benchmarks are conducted on clinical or genomic datasets from high-income settings, and very few have systematically evaluated imputation methods on incomplete social and behavioural datasets from South Asia \cite{ref8}. Second, most studies evaluate imputation quality in isolation --- measuring reconstruction error against artificially masked values --- rather than the downstream impact of imputation choices on unsupervised tasks such as clustering \cite{ref8,ref55}. Third, limited prior work in the newspaper-derived social-data setting has examined how predictor selection during imputation --- whether all available columns or only contextually correlated predictors are used --- propagates into downstream cluster structure, which is one of the dimensions explored in the present study.

\section{Related Work on Clustering in Suicide and Mental Health Research}

Unsupervised clustering has been increasingly applied to mental health and suicide data as a means of identifying latent subgroups within heterogeneous populations; unlike supervised approaches, clustering lets the data reveal structure, an appropriate strategy when no ground-truth class labels exist. At the broadest level, Orchard et al. (2024) \cite{ref34} applied K-Means and hierarchical clustering to a large Canadian community health survey and identified four mental health profiles with meaningfully different healthcare service-use patterns, demonstrating that unsupervised methods can surface clinically and policy-relevant subgroups invisible to traditional regression.

More directly relevant to the suicide context, Xiao et al. (2025) \cite{ref35} published what is described as the first study to apply unsupervised machine learning to a comprehensive set of social determinants of health for suicide risk across U.S. communities; using K-Means on large population-level data, they identified geographically coherent community-level risk clusters --- including one characterised by social isolation and poor healthcare access, and another by economic distress and family instability --- each requiring distinct prevention strategies. The motivations parallel those of the present study, though the U.S. context, structured administrative data, and community-level (rather than individual-level) variables make the methodological challenges substantially different from those of incomplete Bangladeshi newspaper data.

Sanz-Gomez et al. (2025) \cite{ref36} applied agglomerative nesting followed by K-Means to psychological autopsy data from 391 suicide deaths in Spain, identifying three impulsivity-based subprofiles and highlighting the heterogeneity that motivates clustering over aggregate analyses. Wong et al. (2022) \cite{ref37} used cluster analysis on stressor and risk-factor data from student suicides in Hong Kong, identifying four meaningful profiles (school distress, hidden, family and relationship, and numerous issues) with distinct prevention implications. More recent supervised ML on adolescent suicide has been comprehensively synthesised by Liu et al. (2025) \cite{ref53}, whose meta-analysis of 42 studies and 104 ML models confirms that ensemble methods (random forest, XGBoost) currently set the performance ceiling for predictive risk modelling in this age group.

In the Bangladeshi context specifically, prior studies have not applied any form of algorithmic clustering to suicide data. Arafat et al. (2018) \cite{ref4}, one of the most cited newspaper-based analyses of Bangladeshi suicides, explicitly acknowledged that high missingness precluded any systematic subgroup analysis beyond basic cross-tabulation --- a direct statement of the methodological gap the present study addresses. Table~\ref{tbl:summary_of_related_work_on_clu} summarises the most relevant clustering studies in the domain.

\begin{table}[H]
\centering
\caption[Related work on clustering in suicide and mental health]{Summary of related work on clustering in suicide and mental health research}
\label{tbl:summary_of_related_work_on_clu}
\renewcommand{\arraystretch}{1.0}
\small
\begin{tabularx}{\textwidth}{|>{\hsize=0.691\hsize\linewidth=\hsize\raggedright\arraybackslash}X|>{\hsize=0.841\hsize\linewidth=\hsize\raggedright\arraybackslash}X|>{\hsize=0.892\hsize\linewidth=\hsize\raggedright\arraybackslash}X|>{\hsize=1.576\hsize\linewidth=\hsize\raggedright\arraybackslash}X|}
\hline
\textbf{Study} & \textbf{Country / Data} & \textbf{Clustering Approach} & \textbf{Main Finding} \\ \hline
Orchard et al. (2024) \cite{ref34} & Canada (community health survey) & K-Means, hierarchical clustering & Four distinct mental health profiles identified, each associated with different healthcare service-use patterns. \\ \hline
Xiao et al. (2025) \cite{ref35} & USA (social determinants data) & Unsupervised machine learning (K-Means) & First study to use unsupervised ML on social determinants of health for suicide; community-level risk clusters identified. \\ \hline
Sanz-Gomez et al. (2025) \cite{ref36} & Spain (psychological autopsy, 391 cases) & Agglomerative nesting + K-Means & Three impulsivity-based suicide profiles identified; highlights significant heterogeneity among suicide decedents. \\ \hline
Wong et al. (2022) \cite{ref37} & Hong Kong (student suicides) & Cluster analysis on stressors and risk factors & Four student suicide profiles found (school distress, hidden, family/relationship, numerous issues); supports differentiated interventions. \\ \hline
Arafat et al. (2018) \cite{ref4} & Bangladesh (newspaper data) & Manual grouping only & No algorithmic clustering applied; missingness in newspaper data cited as primary barrier to advanced analysis. \\ \hline
\end{tabularx}
\end{table}

A critical observation across the reviewed clustering studies is that, within the works surveyed for this thesis, none systematically varied both the imputation strategy and the clustering algorithm to understand how imputation choices affect downstream cluster quality. Almost all relied on standard internal validity indices --- Silhouette score, Calinski--Harabasz, Davies--Bouldin --- which the recent comprehensive review of cluster validity indices catalogues alongside their known sensitivities, including degraded behaviour under severe cluster-size imbalance \cite{ref54}: when one cluster contains the vast majority of observations and another only a handful of outliers, standard validity indices can still report high scores, making the solution appear better than it is in practice. This limitation is particularly acute in real-world social datasets where natural group imbalance is the norm rather than the exception.

\section{Problem Analysis and Research Gap}

From the foregoing review, the following critical gaps emerge and directly motivate the design of the present study:

\begin{enumerate}[leftmargin=2em,labelindent=0pt,itemsep=0.4em]
  \item \textbf{No large-scale, enriched suicide dataset for Bangladesh.} The largest existing newspaper-based Bangladeshi suicide datasets contain at most a few hundred records with limited attributes and cover single years. No publicly available dataset systematically covers multiple years (2017--2025) or includes socioeconomic, clinical, and environmental attributes alongside basic demographics.
  \item \textbf{No systematic imputation benchmark for this data type.} The international imputation literature is extensive, but no study has systematically benchmarked state-of-the-art imputation methods (MICE, KNN, MissForest) on a large, real-world, newspaper-derived social dataset from South Asia at the extreme missingness levels characteristic of this domain (50--90\% for several key variables).
  \item \textbf{No multi-method clustering study on Bangladeshi suicide data.} No published work has applied multiple clustering algorithms to Bangladeshi suicide records, let alone examined how the choice of imputation method influences the resulting cluster structure. Existing clustering applications in this domain have been conducted exclusively on Western datasets with structured, near-complete data.
  \item \textbf{No evaluation framework accounting for cluster imbalance.} Standard internal clustering validity indices do not penalise solutions in which a single cluster dominates the data or density-based algorithms classify the majority of points as noise. For naturally imbalanced social datasets, a fairness-aware evaluation component is needed to distinguish practically useful clusterings from statistically artefactual ones.
  \item \textbf{No consideration of predictor-selection effects during imputation.} Existing comparative imputation studies rarely examine how the choice of predictor set --- all available columns versus only contextually correlated ones --- affects both imputation quality and downstream clustering outcomes. We are not aware of prior work examining this dimension specifically in the Bangladeshi public-health context.
\end{enumerate}

The present study addresses all five gaps by expanding the publicly available dataset from 760 to 1,617 cases with 9 new attributes covering 2017--2025; implementing a 32-pipeline comparison of three imputation methods under two predictor strategies and four clustering algorithms; introducing a composite scoring formulation that combines five standard validity indices with a study-specific balance penalty term; and providing interpretable, hypothesis-generating cluster descriptions grounded in the Bangladeshi social context. The work also draws methodologically on the broader imputation-benchmark literature: Jarrett et al. (2021) \cite{ref8} concluded that no single imputation method dominates across data types and missingness patterns, which directly motivates the multi-method comparison adopted here. Table~\ref{tbl:comparison_with_prior_works} positions the present study against three representative prior suicide studies across methodological and contextual dimensions.

\begin{table}[H]
\centering
\caption[Comparison of present study with key prior suicide studies]{Comparison of the present study with key prior suicide studies across nine methodological and contextual dimensions}
\label{tbl:comparison_with_prior_works}
\renewcommand{\arraystretch}{1.0}
\small
\begin{tabularx}{\textwidth}{|>{\hsize=0.68\hsize\linewidth=\hsize\raggedright\arraybackslash}X|>{\hsize=0.75\hsize\linewidth=\hsize\raggedright\arraybackslash}X|>{\hsize=0.75\hsize\linewidth=\hsize\raggedright\arraybackslash}X|>{\hsize=0.85\hsize\linewidth=\hsize\raggedright\arraybackslash}X|>{\hsize=1.97\hsize\linewidth=\hsize\raggedright\arraybackslash}X|}
\hline
\textbf{Dimension} & \textbf{Arafat et al. 2018 \cite{ref4}} & \textbf{Xiao et al. 2025 \cite{ref35}} & \textbf{Sanz-Gomez et al. 2025 \cite{ref36}} & \textbf{This Study} \\ \hline
Country / Context & Bangladesh & USA & Spain & Bangladesh \\ \hline
Data source & Newspaper & Administrative & Clinical autopsy & Newspaper \\ \hline
Dataset size & \textasciitilde{}360 cases & Large population & 391 cases & 1,617 cases \\ \hline
Missing data handling & None / complete case & Not applicable & Not applicable & MICE + KNN + MissForest $\times$ 2 strategies \\ \hline
Clustering applied & None & K-Means only & Agglomerative + K-Means & 4 algorithms $\times$ 8 datasets = 32 pipelines \\ \hline
Evaluation metric & Manual grouping & Standard metrics & Standard metrics & 5 standard + study-specific balance penalty \\ \hline
Multi-year coverage & 2009--2018 & Multiple years & Single study & 2017--2025 \\ \hline
Reproducible pipeline & No & Partial & No & Yes (12 notebooks, open source) \\ \hline
Study-specific scoring component & No & No & No & Yes (balance penalty term) \\ \hline
\end{tabularx}
\end{table}

Read together with the gaps above and the per-cell comparison in Table~\ref{tbl:comparison_with_prior_works}, the position of this study against the closest prior work can be stated precisely. Arafat et al. (2018) \cite{ref4} is the most directly comparable Bangladeshi precedent in terms of data source and country context, but it does not apply any algorithmic clustering and explicitly attributes that limitation to the unresolved missingness problem in newspaper-derived data --- exactly the problem the present study is built to address. Xiao et al. (2025) \cite{ref35} and Sanz-Gomez et al. (2025) \cite{ref36} demonstrate that algorithmic subgroup discovery can yield interpretable, intervention-relevant clusters in suicide and mental-health populations, but they work with administrative or clinical-autopsy data that does not exhibit the multi-source, $>$50\% missingness profile of Bangladeshi newspaper records, and each applies a single clustering family rather than comparing across imputation and clustering choices jointly. The recent systematic review of machine-learning approaches to adolescent suicide risk \cite{ref53}, taken together with the comprehensive review of cluster validity indices \cite{ref54}, likewise treats imputation and clustering as separable methodological choices, and the surveyed work does not report side-by-side comparisons of multiple imputation strategies under matched clustering pipelines. The distinctive contribution of the present work, viewed against this prior literature, is therefore not a new clustering algorithm but a domain-specific methodological package: (i) the largest curated Bangladeshi newspaper-derived suicide dataset we are aware of at the time of writing (1,617 records, expanded from $\sim$760), (ii) the first joint comparison of three imputation strategies under two predictor regimes against four clustering algorithms on this data type (32 end-to-end pipelines), and (iii) a balance-aware composite evaluation framework that explicitly demotes the degenerate-but-tight partitions which single-index validation rewards on imbalanced social data. The substantive findings (for example, the source-cohort confound reported in Chapter~4) are themselves a methodological contribution of this design --- they are recoverable precisely because the comparative framework was built to surface, rather than hide, structural artefacts of multi-source data.

\section{Summary}

Suicide in Bangladesh is a serious and growing public-health concern characterised by predominantly hanging-method deaths, over-representation of young females and students, sharp rural--urban disparities, and high-burden districts such as Jhenaidah and Chuadanga, with case-level data almost entirely confined to newspaper reports that suffer extreme, structurally non-random missingness. The international imputation literature consistently favours MissForest on mixed-type data, but no benchmark has been conducted in the specific context of newspaper-based South Asian social data, and no prior study on this data type has jointly compared multiple imputation strategies against multiple clustering algorithms or used a balance-aware metric.

Each of the five gaps identified in \S2.6 maps to a specific section of the methodology that follows. The dataset gap (gap~1) is addressed by the case-collection and column-cleaning steps of \S3.2, which together produce the 1{,}617-record corpus on which everything downstream depends. The imputation-benchmark gap (gap~2) is addressed by the three-method comparison in \S3.4, where MICE, KNN, and MissForest are run under standardised parameter and convergence settings. The multi-method clustering gap (gap~3) is addressed by \S3.5, which applies four algorithms drawn from distinct paradigms to all eight input datasets and yields the 32-pipeline grid. The evaluation-metric gap (gap~4) is addressed by \S3.6.3, which defines the study-specific balance penalty that demotes degenerate-but-tight partitions on naturally imbalanced data. The predictor-selection gap (gap~5) is addressed by the context-aware versus no-context comparison built into the imputation stage of \S3.4.1.1. Chapter~3 sets out the integrated framework that closes these gaps and produces the experimental results analysed in Chapter~4.
